% Chapter 14, generated by make_chapters.py from papers/Ch14_Bio_10_Pangenome_Sheaf/anccft23.tex.
% Source paper: Bio 10
\chapter[The recombination group of a pan-genome]{The pan-genome sheaf sees nothing: the recombination group of a pan-genome graph, and its torsion}\label{chap:ch14}
\noindent{\small\emph{Paper Bio~10 of the series. Verification script and transcript: \texttt{papers/Ch14\_Bio\_10\_Pangenome\_Sheaf/}.}}\par\medskip
\begin{chaptersummary}
A pan-genome graph is the pooled de Bruijn multigraph $G$ of strains $H_1,\dots,H_p$ together
with the strain walks, so the decomposition of the pooled flow into strains, which the
phasing note had to count, is here given. The natural cellular sheaf on $G$ --- stalk the
free abelian group on the strains through a vertex, restriction maps the projections --- is
the direct sum of the constant sheaves on the strain subgraphs $G_i$: its cohomology is
$\Z^p$ and $\bigoplus_i H^1(G_i)$, and the integer cokernel of its sheaf Laplacian is
$\bigoplus_i(\Z\oplus\K(|G_i|))$. It contains no pan-genomic information. What does is the
\emph{recombination group} $\RecGrp=Z_1(G)/\sum_iZ_1(G_i)$, the circulations of the pooled
graph modulo those that live within single strains: the cycles that exist only by switching
strain. The Mayer--Vietoris spectral sequence of the cover $\{G_i\}$ identifies $\RecGrp$ with the
first homology of the nerve of intersection components, $C_q=\bigoplus_{|\sigma|=q+1}
H_0(G_\sigma)$. For two strains $\RecGrp$ is free of rank $c(G_1\cap G_2)-1$ --- the number of
shared blocks minus one, which on a repeat-free reference with $b$ bubbles is $b-1$ and on
$\phi$X174, whose two-copy repeat glues shared blocks together, is smaller. For three or more
strains $\RecGrp$ can have torsion: six synthetic strains whose intersection pattern is the
six-vertex triangulation of the projective plane have $\RecGrp=\Z/2$, while the octahedron gives
$0$ and the seven-vertex torus $\Z^2$. A torsion class is a mosaic cycle twice of which, but
not itself, is a sum of within-strain cycles; no single strain and no pair can see it. Six
$\phi$X174 isolates from GenBank, strand- and origin-normalised, give $\RecGrp=\Z^{38}$ with all
$63$ intersections non-empty and no torsion. Every group is computed twice, by a Smith form of
the inclusion of the strain cycle lattices and from the nerve, and the two agree in all $28$
pools. On complete \emph{E.\ coli} chromosomes --- where the cycle lattice is far out of
reach but the nerve is not --- $\RecGrp$ is free, and freely so at every prime at once: the
\v Cech boundary $d_2$ reduces to the identity under $\pm1$ pivots, which certifies that all
its elementary divisors are $1$. Torsion needs a tight triangulated surface, and a real
pan-genome nerve, whose strains share their core in tens of thousands of separate blocks, is
the opposite shape. The certificate holds on every pool from two to eight chromosomes, up to
$\RecGrp=\Z^{117\,848}$ on a nerve with $1.7$ million two-cells, and does so even where $d_2$ has
a large kernel; the rank of $\RecGrp$ saturates as strains are added, measuring the diversity of
the set rather than its size.
\end{chaptersummary}

\medskip
\noindent\textbf{Keywords.} pan-genome graph, variation graph, cellular sheaf, sheaf
Laplacian, cycle space, Mayer--Vietoris spectral sequence, nerve, recombination, torsion.

\noindent\textbf{MSC 2020.} 05C20, 05C50, 55N30, 55T99, 92D20.

\section{What was asked, and what is true}

The plan for this volume listed a note on ``cellular sheaves on pan-genome graphs'' and ranked
it last, with the warning that the sheaf might turn out to be a relabelling of what the
phasing note already had. This note carries out the construction and reports that the
warning was right in one direction and wrong in another. The cellular sheaf that a pan-genome
graph carries naturally is trivial in a precise sense (Proposition~\ref{ch14:prop:sheaf}); but
the question the sheaf was meant to ask --- what does the pooled graph know that no single
strain knows --- has a clean answer that is not a relabelling of anything, and it comes with
an integer invariant that can be torsion.

A pan-genome graph is a sequence graph together with the walks of the individual genomes
through it \cite{GSN18,EPN20}. In the de Bruijn setting of these notes, that is the pooled
multigraph $G=G_k(H_1{+}\dots{+}H_p)$ of Chapter~\ref{chap:ch13} with the decomposition
$m=\sum_i\spec(H_i)$ \emph{known}. The phasing note counted the decompositions; here there is
one, and the object of interest is the relation between the strain subgraphs $G_i$ and their
union. The cellular sheaf $\sheaf$ that records, at each vertex and arc, which strains pass
through it is the sheaf-theoretic encoding of that relation, and by the spectral theory of
cellular sheaves \cite{HG19} one would compute its cohomology and its sheaf Laplacian.
Proposition~\ref{ch14:prop:sheaf} says both decompose strain by strain: $\sheaf$ is the direct sum
of the constant sheaves on the $G_i$, and everything one computes from it is a list of
per-strain quantities. The sheaf sees no interaction because the projections that define it
never mix strains.

The interaction lives one step further along. The circulation lattice $Z_1(G)$ of the pooled
graph contains the circulation lattices $Z_1(G_i)$ of the strains, and the quotient
\[
  \RecGrp(H_1,\dots,H_p;k)=Z_1(G)\Big/\sum_{i=1}^pZ_1(G_i)
\]
is the group of cycles of the pan-genome that cannot be assembled from cycles of individual
genomes --- the cycles a walk can only close by switching from one strain to another at a
shared $k$-mer, which is to say the recombination cycles. Section~\ref{ch14:sec:R} computes it.
The Mayer--Vietoris spectral sequence of the cover of $G$ by the $G_i$ identifies $\RecGrp$ with
the first homology of a small chain complex built from the connected components of the
intersections $G_{i_0}\cap\dots\cap G_{i_q}$ (Theorem~\ref{ch14:thm:mv}). Two consequences
follow. For two strains, $\RecGrp$ is free of rank one less than the number of shared blocks
(Corollary~\ref{ch14:cor:two}). For three or more, $\RecGrp$ can have torsion, and
Proposition~\ref{ch14:prop:rp2} exhibits six strains with $\RecGrp=\Z/2$. Section~\ref{ch14:sec:data}
verifies all of it, and computes $\RecGrp$ for six deposited $\phi$X174 genomes.

Throughout, strains are circular sequences over $\{A,C,G,T\}$, $k$ is fixed, and the pooled
graph is unitig-compacted with the branch criterion of Chapter~\ref{chap:ch13} (a $k$-mer is a vertex if
it has two distinct successors or two distinct predecessors in the pool); a \emph{content}
is a unitig string, $A$ the set of contents, and $Z_1(G)\subset\Z^A$ the circulation lattice.
All computations are exact; the script \texttt{pangenome\_sheaf.py} reproduces every number
below and prints a five-line pass/fail battery.

\section{The pan-genome sheaf}\label{ch14:sec:sheaf}

\begin{definition}\label{ch14:def:sheaf}
Let $G_i\subseteq G$ be the subgraph of contents and vertices used by strain $i$, and for a
vertex or content $x$ let $S(x)=\{i:x\in G_i\}$. The \emph{pan-genome sheaf} $\sheaf$ has
stalks $\sheaf(x)=\Z^{S(x)}$ and, for a content $a$ with endpoint $v$, the restriction map
$\sheaf(v)\to\sheaf(a)$ the coordinate projection $\Z^{S(v)}\to\Z^{S(a)}$ (note
$S(a)\subseteq S(v)$). Its coboundary $\delta\colon\bigoplus_v\sheaf(v)\to\bigoplus_a\sheaf(a)$
is $(\delta x)_a=x_{\mathrm{head}(a)}|_{S(a)}-x_{\mathrm{tail}(a)}|_{S(a)}$, and its sheaf
Laplacian is $L_\sheaf=\delta^{\mathsf T}\delta$.
\end{definition}

\begin{proposition}\label{ch14:prop:sheaf}
$\sheaf\cong\bigoplus_{i=1}^p\Z_{G_i}$, the direct sum of the constant sheaves $\Z$ supported
on the strain subgraphs. Consequently $H^0(G;\sheaf)\cong\Z^p$,
$H^1(G;\sheaf)\cong\bigoplus_iH^1(G_i;\Z)$, $L_\sheaf$ is block diagonal with blocks the
graph Laplacians of the undirected shadows $|G_i|$, and
\[
  \coker L_\sheaf\;\cong\;\bigoplus_{i=1}^p\big(\Z\oplus\K(|G_i|)\big),
\]
$\K(|G_i|)$ the critical group of the undirected multigraph $|G_i|$.
\end{proposition}

\begin{proof}
The coordinate $i$ of every stalk is present exactly on the cells of $G_i$, and the
restriction maps are the identity on that coordinate wherever both cells carry it and zero
otherwise; that is the constant sheaf $\Z$ on the subcomplex $G_i$, extended by zero. Direct
sums of sheaves have direct sums of cochain complexes, so cohomology and Laplacian decompose;
the cochain complex of $\Z_{G_i}$ is the cellular cochain complex of the graph $G_i$, whose
$H^0$ is $\Z$ for connected $G_i$ (a strain walk is connected) and whose Laplacian is the
graph Laplacian of $|G_i|$, with cokernel $\Z\oplus\K(|G_i|)$.
\end{proof}

Check (F) computes $\coker L_\sheaf$ for a diploid with three substitutions and for three
strains, obtaining $\Z/38\oplus\Z/38\oplus\Z^2$ and $(\Z/84)^3\oplus\Z^3$ respectively, and
confirms the block decomposition by a Smith form of the assembled Laplacian. The numbers
$38$ and $84$ are the critical groups of the strains' own shadows; they would be the same if
the strains were sequenced separately. There is nothing pan-genomic in the sheaf, and no
choice of stalk dimensions rescues this as long as the restriction maps are projections: a
sheaf whose restriction maps never mix coordinates is a direct sum of coordinate sheaves.

\section{The recombination group}\label{ch14:sec:R}

\begin{definition}
$\RecGrp=Z_1(G)/\sum_iZ_1(G_i)$, where $Z_1(G_i)=Z_1(G)\cap\Z^{A_i}$ is the circulation
lattice of the subgraph $G_i$. For a subset $\sigma\subseteq\{1,\dots,p\}$ write
$G_\sigma=\bigcap_{i\in\sigma}G_i$ and let $C_q=\bigoplus_{|\sigma|=q+1}H_0(G_\sigma)$, the
free abelian group on the connected components of the $(q{+}1)$-fold intersections, with the
\v Cech differential $\partial(\sigma,c)=\sum_j(-1)^j(\sigma\setminus\sigma_j,\,c')$, $c'$
the component of $G_{\sigma\setminus\sigma_j}$ containing $c$. We call $(C_\bullet,\partial)$
the \emph{nerve of components} $\Nerve$.
\end{definition}

\begin{theorem}\label{ch14:thm:mv}
$\RecGrp\cong H_1(\Nerve)$.
\end{theorem}

\begin{proof}
The subgraphs $G_i$ are subcomplexes covering the one-dimensional CW complex $G$, so the
Mayer--Vietoris spectral sequence of the cover \cite[VII.4]{Bro82} applies:
$E^1_{q,r}=\bigoplus_{|\sigma|=q+1}H_r(G_\sigma)\Rightarrow H_{q+r}(G)$, with $d^1$ the
\v Cech differential. Since intersections of graphs are graphs, $H_r(G_\sigma)=0$ for
$r\ge2$, and the $E^2$ page has two rows, $E^2_{q,0}=H_q(\Nerve)$ and
$E^2_{q,1}=H_q\big(\bigoplus_\sigma H_1(G_\sigma)\big)$. The only differentials that can
touch total degree one are $d^2\colon E^2_{2,0}\to E^2_{0,1}$ and the vanishing
$d^2\colon E^2_{1,0}\to E^2_{-1,1}$, so $E^\infty_{1,0}=E^2_{1,0}=H_1(\Nerve)$ and
$E^\infty_{0,1}=E^2_{0,1}/d^2(E^2_{2,0})$. The filtration of $H_1(G)$ has $F_0=E^\infty_{0,1}$
equal to the image of $\bigoplus_iH_1(G_i)\to H_1(G)$ (the edge homomorphism) and
$H_1(G)/F_0\cong E^\infty_{1,0}$. For a graph $H_1=Z_1$, and the image of
$\bigoplus_iZ_1(G_i)$ in $Z_1(G)$ is $\sum_iZ_1(G_i)$.
\end{proof}

\begin{corollary}[Two strains]\label{ch14:cor:two}
For $p=2$ with $G_1,G_2$ connected and $G_1\cap G_2\ne\emptyset$, $\RecGrp$ is free of rank
$c(G_1\cap G_2)-1$, the number of connected components of the shared subgraph minus one.
\end{corollary}

\begin{proof}
$\Nerve$ has $C_0=\Z^2$, $C_1=\Z^{c(G_{12})}$ and $C_2=0$; $\partial$ sends every component to
$e_2-e_1$, so $H_1(\Nerve)=\ker\partial$ is the sublattice of coefficient sum zero.
\end{proof}

On a repeat-free reference with $b$ substitution sites in distinct $(k{+}1)$-windows, the
shared subgraph is the $b$ stretches between the bubbles and $\RecGrp\cong\Z^{b-1}$. When the
reference has a repeat the shared stretches that contain copies of the same repeated $k$-mer
are joined through it and $c(G_1\cap G_2)$ drops: for $\phi$X174 at $k=12$, whose two
repeated $12$-mers $a,b$ each occur twice, five substitutions leave three shared components and
$\RecGrp=\Z^2$, three substitutions leave two and $\RecGrp=\Z$ (Table~\ref{ch14:tab:two}). A repeat
inside a strain thus reduces the number of independent recombination cycles between strains:
the cycle through the repeat is already a within-strain cycle.

\begin{proposition}[Torsion]\label{ch14:prop:rp2}
Let $T$ be a triangulated closed surface with vertex set $\{1,\dots,p\}$, and construct
strains as follows: a random segment $F_\tau$ for each face $\tau$ and $E_\varepsilon$ for
each edge $\varepsilon$; strain $i$ reads, for each edge $\varepsilon=ij$ at $i$ with faces
$\tau,\tau'$, the word $F_\tau E_\varepsilon F_{\tau'}\,S\,F_{\tau'}E_\varepsilon F_\tau\,S'$
with unique spacers $S,S'$. Then for $k$ larger than any accidental shared $k$-mer the nerve
of components is the simplicial chain complex of $T$, and $\RecGrp\cong H_1(T;\Z)$. In
particular the six-vertex triangulation of the projective plane gives $\RecGrp=\Z/2$, the
octahedron gives $0$, and the seven-vertex torus gives $\Z^2$.
\end{proposition}

\begin{proof}
$G_i\cap G_j$ for an edge $ij$ consists of $E_{ij}$, the two face segments at $ij$, and the
four junction arcs between them, which both strains contain because each reads the two
junction directions; it is connected. For a face $ijk$, $G_i\cap G_j\cap G_k=F_{ijk}$, a
segment; for a non-face triple and for all larger subsets the intersection is empty, since
every segment is shared by the vertices of exactly one simplex and every junction arc by
exactly two strains. So $C_\bullet$ is the chain complex of $T$ and Theorem~\ref{ch14:thm:mv}
applies. $H_1(\RP;\Z)=\Z/2$, $H_1(S^2)=0$, $H_1(T^2)=\Z^2$.
\end{proof}

The torsion class is concrete. In the projective-plane pool there is a circulation $z$ of the
pooled graph such that $2z$ is a sum of cycles each lying within a single strain but $z$
itself is not; $z$ is a mosaic that closes up only by using every strain, and its double
decomposes. No pair of strains detects it, because for pairs $\RecGrp$ is free
(Corollary~\ref{ch14:cor:two}); no strain detects it, because within a strain there is nothing to
detect. It is an invariant of the pan-genome as a whole.

\section{Verification and data}\label{ch14:sec:data}

Every group is computed twice: as the cokernel, via Smith normal form, of the matrix whose
columns are fundamental cycles of the subgraphs $G_i$ expressed in $\Z^A$ (torsion of
$\Z^A/\sum Z_1(G_i)$ equals that of $\RecGrp$ because $Z_1(G)$ is saturated in $\Z^A$), and as
$H_1(\Nerve)$ from the component structure of the intersections. The two agree in every one
of the $28$ pools, which is check (T) and part of the others.

\begin{table}[t]
\centering
\small
\begin{tabular}{lrrrrrl}
\toprule
pool ($k=12$) & $|V|$ & $|A|$ & $\rank Z_1(G)$ & $\rank\sum Z_1(G_i)$ & $c(G_1\cap G_2)$ & $\RecGrp$\\
\midrule
$\phi$X174, 1 substitution          &  4 &  7 & 4 & 4 & 1 & $0$\\
$\phi$X174, 2 substitutions         &  6 & 10 & 5 & 4 & 2 & $\Z$\\
$\phi$X174, 3 substitutions         &  8 & 13 & 6 & 5 & 2 & $\Z$\\
$\phi$X174, 5 substitutions         & 12 & 19 & 8 & 6 & 3 & $\Z^2$\\
$\phi$X174, sites $600,606$ linked $+1$ & 6 & 10 & 5 & 4 & 2 & $\Z$\\
repeat-free reference, 2 sites      &  4 &  6 & 3 & 2 & 2 & $\Z$\\
repeat-free reference, 3 sites      &  6 &  9 & 4 & 2 & 3 & $\Z^2$\\
repeat-free reference, 4 sites      &  8 & 12 & 5 & 2 & 4 & $\Z^3$\\
\bottomrule
\end{tabular}
\caption{Two strains. $\RecGrp$ is free of rank $c(G_1\cap G_2)-1$ in every row; on the
repeat-free reference (a random $3000$-mer) that is the number of bubbles minus one, on
$\phi$X174 the repeat glues shared stretches and lowers it. Check (B).}
\label{ch14:tab:two}
\end{table}

\begin{table}[t]
\centering
\small
\begin{tabular}{lrrrrrl}
\toprule
surface pool ($k=16$) & strains & $|V|$ & $|A|$ & $\rank Z_1(G)$ & $\rank\sum Z_1(G_i)$ & $\RecGrp=H_1(\Nerve)$\\
\midrule
octahedron ($S^2$)               & 6 & 114 & 190 &  77 &  77 & $0$\\
seven-vertex torus ($T^2$)       & 7 & 186 & 319 & 134 & 132 & $\Z^2$\\
hemi-icosahedron ($\RP$)         & 6 & 138 & 233 &  96 &  96 & $\Z/2$\\
\bottomrule
\end{tabular}
\caption{Proposition~\ref{ch14:prop:rp2}. In each pool the nerve of components is verified to be
exactly the designed triangulation (one component per simplex, none elsewhere). At $k=12$
the $\approx10^4$ random $12$-mers of a pool produce accidental shared $k$-mers and spurious
components; $k=16$ removes them. Check (S).}
\label{ch14:tab:surfaces}
\end{table}

\begin{table}[t]
\centering
\small
\begin{tabular}{llrrrrrr}
\toprule
 & accession & \multicolumn{6}{c}{substitutions after normalisation}\\
\midrule
1 & LC786485.1 & 0 & 5 & 27 & 26 & 24 & 33\\
2 & NC\_001422.1 & 5 & 0 & 26 & 25 & 23 & 32\\
3 & MH378443.1 & 27 & 26 & 0 & 19 & 24 & 48\\
4 & MH378442.1 & 26 & 25 & 19 & 0 & 22 & 47\\
5 & MH378441.1 & 24 & 23 & 24 & 22 & 0 & 45\\
6 & JQ764988.2 & 33 & 32 & 48 & 47 & 45 & 0\\
\bottomrule
\end{tabular}
\caption{Six $\phi$X174 genomes of length $5386$ from GenBank, brought to the strand and
origin of NC\_001422.1 (one was deposited rotated). Pooled at $k=12$: $|V|=104$, $|A|=158$,
$\rank Z_1(G)=55$, $\rank\sum Z_1(G_i)=17$, $\RecGrp=\Z^{38}$, no torsion; all $63$
intersections are non-empty, with $4$ to $27$ components for pairs and $41$ for the full
intersection. Check (D).}
\label{ch14:tab:iso}
\end{table}

\paragraph{The checks.} In \texttt{pangenome\_sheaf.py}: (F) the sheaf Laplacian's cokernel is
the direct sum of the strains' $\Z\oplus\K(|G_i|)$; (B) two strains, eight pools, $\RecGrp$ free
of rank $c(G_1\cap G_2)-1$; (S) the three surface pools, with the designed nerve; (T) twelve
random three-strain pools from a shared segment library, Smith form and nerve agreeing, no
torsion; (D) the six isolates. All pass in seven seconds. In \texttt{real\_pangenome.py},
for Section~\ref{ch14:sec:real}: (A) the hashed branch criterion against brute force on random
$20$\,kbp fragment pools; (V) the $\mathbb F_p$ detector against the three surface pools,
where it returns the $\Z/2$ of the projective plane; (E) the complete chromosomes of
Table~\ref{ch14:tab:real}.

\section{A real pan-genome}\label{ch14:sec:real}

Theorem~\ref{ch14:thm:mv} replaces a lattice quotient by the homology of the nerve, and the nerve
needs only the \emph{connected components} of the shared subgraphs. Components are cheap:
the recombination group is therefore computable on complete bacterial chromosomes, where
$\RecGrp$ itself, as a quotient of a cycle lattice of rank $10^4$ or more, is not. This section
carries that out.

\paragraph{The pool.} Complete \emph{Escherichia coli} chromosomes were taken from GenBank
and brought to the strand of K-12 MG1655 by $k$-mer overlap; a circular genome's $k$-mer set
is rotation invariant, so only the strand has to be fixed. The pooled graph is compacted by
the criterion of Chapter~\ref{chap:ch13}: a $k$-mer is a vertex when it has two distinct successors or
two distinct predecessors \emph{in the pool}. Branch vertices are found by hashing on the
pairs (k-mer, next base) and (k-mer, previous base) and then re-tested on the exact strings,
so a collision can only propose a candidate that the string pass discards; check (A) audits
the whole criterion against brute force on random $20$\,kbp fragment pools.

\paragraph{A certificate, not a sample.} Torsion is a statement about every prime, and one
cannot test them all. Two devices avoid having to.

\begin{proposition}[freeness certificate]\label{ch14:prop:cert}
Suppose the \v Cech boundary $d_2\colon C_2\to C_1$ can be reduced by integer elimination in
which every pivot is $\pm1$. Then all its elementary divisors are $1$, $\im d_2$ is a direct
summand of $C_1$, and $\RecGrp=\ker d_1/\im d_2$ is free.
\end{proposition}

\begin{proof}
A $\pm1$ pivot is a unimodular row and column operation, so a complete reduction exhibits
$d_2$ as $\binom{I}{0}$ up to unimodular change of basis; its elementary divisors are all
$1$ and $C_1/\im d_2$ is free. A finitely generated submodule of a finitely generated free
$\Z$-module is free, and $\RecGrp=\ker d_1/\im d_2\subseteq C_1/\im d_2$.
\end{proof}

\begin{remark}[a torsion detector without a Smith form]\label{ch14:rem:detector}
When the reduction does not complete, the prime-by-prime test is still cheap. $H_0(\Nerve)$
is free, so the universal coefficient theorem gives
$\dim_{\mathbb F_p}H_1(\Nerve;\mathbb F_p)=\rank\RecGrp+t_p(\RecGrp)$ with $t_p$ the $p$-torsion
rank: any excess of the $\mathbb F_p$ dimension over the rational rank is $p$-torsion, and
both are ranks of sparse matrices. Check (V) runs this detector on the three surface pools of
Table~\ref{ch14:tab:surfaces}, where the answer is known exactly, and it does return the $\Z/2$ of
the projective plane --- so a null result on real data is a result and not a blind spot.
\end{remark}

\begin{table}[t]
\centering
\small
\begin{tabular}{lrrrrrll}
\toprule
strains ($k=31$) & branch vtx & contents & $C_1$ & $C_2$ & $\rank d_2$ & $\RecGrp$ & certified\\
\midrule
MG1655, DH10B & 1\,509 & 2\,343 & 86 & 0 & 0 & $\Z^{85}$ & yes\\
$+$ O157:H7 Sakai & 63\,245 & 95\,168 & 58\,454 & 28\,820 & 28\,820 & $\Z^{29\,632}$ & yes\\
$+$ CFT073 & 137\,375 & 207\,028 & 166\,011 & 134\,401 & 99\,561 & $\Z^{66\,447}$ & yes\\
$+$ IAI39 & 180\,256 & 271\,937 & 260\,369 & 277\,947 & 172\,507 & $\Z^{87\,858}$ & yes\\
$+$ S88 & 214\,033 & 323\,073 & 409\,007 & 588\,230 & 304\,507 & $\Z^{104\,495}$ & yes\\
$+$ 536 & 237\,525 & 358\,660 & 575\,831 & 1\,053\,725 & 459\,769 & $\Z^{116\,056}$ & yes\\
$+$ UTI89 & 241\,592 & 364\,819 & 744\,598 & 1\,683\,865 & 626\,743 & $\Z^{117\,848}$ & yes\\
\bottomrule
\end{tabular}
\caption{The recombination group of eight complete \emph{E.\ coli} chromosomes, about
$40$\,Mbp, each row adding one strain to the one above. In every row the unit-pivot reduction
of $d_2$ completes, so by Proposition~\ref{ch14:prop:cert} $\RecGrp$ is free at every prime, and the
rank, being $\dim C_1-\rank d_1-\rank d_2$ with both ranks computed by unimodular integer
elimination, is exact. The two K-12 strains share $86$ blocks and give $\rank\RecGrp=86-1$,
which is Corollary~\ref{ch14:cor:two} on real data. From the fourth strain on $d_2$ is far from
injective --- at eight strains $626\,743$ of $1\,683\,865$ --- which costs the certificate
nothing: a kernel makes columns vanish during the reduction, it does not stop the surviving
pivots being units. Check (E).}
\label{ch14:tab:real}
\end{table}

\paragraph{The answer.} $\RecGrp$ is free, with a certificate, on every pool from two to eight
complete chromosomes (Table~\ref{ch14:tab:real}). This settles the first open question of an
earlier version of this note for those strains, and the reason is visible in the same table. A
torsion class needs a tight two-dimensional cycle: the projective-plane pool of
Proposition~\ref{ch14:prop:rp2} has exactly one component per simplex, so its nerve \emph{is} a
triangulated surface. A real pan-genome nerve is the opposite shape. Each pair of divergent
strains shares its core genome in tens of thousands of blocks --- $86$ for two K-12 strains,
$28\,817$ and $29\,551$ for K-12 against O157:H7 --- and each triple intersection is more
fragmented still, so the nerve is wide and sparse, every elementary divisor of $d_2$ is $1$,
and $H_1$ is a large free group. The synthetic construction of
Section~\ref{ch14:sec:R} was necessary precisely because real data does not produce that shape,
and the honest reading of Proposition~\ref{ch14:prop:rp2} is that torsion is realisable in
principle rather than observed in \emph{E.\ coli}.

\paragraph{The rank saturates.} One thing in Table~\ref{ch14:tab:real} is not about torsion. The
successive increments of $\rank\RecGrp$ are $29\,547$, $36\,815$, $21\,411$, $16\,637$,
$11\,561$, $1\,792$: after the third strain each new chromosome contributes less than the
one before, and the eighth contributes an order of magnitude less than the fourth, on a
nerve whose two-cells meanwhile grew to $1.7$ million. The pooled graph agrees --- the last
strain adds $4\,067$ branch vertices where the fifth added $42\,881$. The ordering is not
neutral: the last four additions (CFT073, S88, 536, UTI89, with IAI39 between them) are
extraintestinal pathogenic strains, so each arrives with most of its mosaic structure
already represented, and the two K-12 strains at the top are a pair, not a sample. Read as
an observation on these eight genomes rather than a law, it says that $\rank\RecGrp$ is a
measure of the \emph{diversity} of a strain set and not of its size, which is what a count
of independent strain-switching cycles ought to be.

\begin{remark}[a failed certificate is a statement about the elimination]\label{ch14:rem:order}
Proposition~\ref{ch14:prop:cert} is one-sided: it certifies freeness when the reduction completes
and says nothing when it does not, because $d_2$ may have all its elementary divisors equal
to one with no $\pm1$ pivot surviving to the end of one particular elimination. That is not a
theoretical caveat. An earlier version of this computation chose each pivot by scanning the
$64$ shortest live columns, and stalled at four strains with $70\,374$ columns left and at
five with $190\,814$, from which we could conclude only --- via
Remark~\ref{ch14:rem:detector} --- that there is no $2$- and no $3$-torsion. Taking instead the
globally shortest column, and within it the row meeting fewest columns, produces less fill-in,
the entries never leave $\{0,\pm1\}$, and every column is eliminated: those two rows of
Table~\ref{ch14:tab:real} became certificates, and the three below them became feasible at all.
The ranks $99\,561$ and $172\,507$ are the same
under both orderings, which is what says the two runs compute the same map. Two consequences
for the arithmetic. Because the unit-pivot phase is a unimodular column reduction over $\Z$ it
is valid modulo every prime with its pivots still units, so when it \emph{does} stall,
$\rank_{\mathbb F_p}d_2$ is the number of pivots plus the rank of the residual block alone;
and the ordering is worth more than the arithmetic, the five-strain reduction taking $4$
seconds where the scan took $14\,080$.
\end{remark}

\paragraph{What the single-strand convention costs here.} Of the eight strains, IAI39
matched the reference on $2331$ of its sampled $31$-mers forward against $1234$ reverse,
where every other strain ran between $6835$ against $234$ and $2970$ against $109$. That is a
large inversion, and to the
single-strand pooled graph an inverted segment is strain-specific sequence: real homology is
lost and the nerve is cut into more components than it should be. The rank of $\RecGrp$ reported
here is therefore an overcount for divergent strains, in a way a bidirected pooled graph
would repair. The freeness is not affected, being a statement about the shape of the nerve.

\section{Scope, provenance, and what is open}

\paragraph{Scope.} Proposition~\ref{ch14:prop:sheaf}, Theorem~\ref{ch14:thm:mv},
Corollary~\ref{ch14:cor:two} and Proposition~\ref{ch14:prop:cert} hold for arbitrary strains;
Proposition~\ref{ch14:prop:rp2} is a construction; Section~\ref{ch14:sec:real} is a computation on
named genomes and claims nothing beyond them.
The pooled graph is single-stranded, as in Chapter~\ref{chap:ch13}; the reverse-complement
version would pool $D_k$'s and is not attempted. ``Recombination'' is used in the graph sense
--- a cycle that switches strain --- and no claim is made about the biological mechanism that
produced any particular mosaic.

\paragraph{Provenance.} Cellular sheaves on graphs and their Laplacians are \cite{HG19};
pan-genome and variation graphs are \cite{GSN18,EPN20}, and the bubble/snarl decomposition of
such graphs, which is the combinatorial shadow of the components of pairwise intersections,
is \cite{PNE18}. The Mayer--Vietoris spectral sequence for a cover by subcomplexes is
standard, e.g.\ \cite[VII.4]{Bro82}; Theorem~\ref{ch14:thm:mv} is its two-row case read off for
graphs. We are not aware of prior work defining the quotient $Z_1(G)/\sum Z_1(G_i)$ for a
pan-genome graph or observing that it can have torsion; a targeted search found none.

\paragraph{Open.} \emph{First}, torsion in a real pan-genome. Section~\ref{ch14:sec:real} answers
this for the \emph{E.\ coli} chromosomes pooled there --- $\RecGrp$ is free, with a certificate
--- and explains why: the nerve of a real pan-genome is wide and sparse, nothing like the
triangulated surface a torsion class needs. What is open is whether any real pan-genome has
one. A plausible place to look is a set of strains whose shared blocks are \emph{few and
large}, so that the nerve is small enough to close up: closely related strains with a
handful of recombination breakpoints, rather than the divergent ones used here. \emph{Second},
the relation of $H_1(\Nerve)$ to the snarl tree of \cite{PNE18}, which organises the pairwise
components hierarchically. \emph{Third}, a bidirected pooled graph, so that an inversion is
homology rather than strain-specific sequence; by the last paragraph of Section~\ref{ch14:sec:real}
this changes the rank of $\RecGrp$ for divergent strains, and Chapter~\ref{chap:ch06}'s anti-automorphism theorem
says the repair cannot be a quotient by the strand involution. \emph{Fourth}, a sheaf that
does see interaction: the restriction maps would
have to mix strains, and the natural candidate is the sheaf of local decompositions of the
pooled flow, whose stalk at a bubble is the homogeneous space of Chapter~\ref{chap:ch13}; it is not a
sheaf of abelian groups.
